\ifx\allfiles\undefined
\documentclass[12pt, a4paper, oneside, UTF8]{ctexbook}
\def\path{../config}
\input{\path/package.tex}

% 定理环境设置
\input{\path/theorem1_zh.tex}

\input{\path/custom.tex}

% 修改参数改变封面样式，0 默认原始封面、内置其他1、2、3种封面样式
\def\myIndex{3}


\ifnum\myIndex>0
    \input{\path/cover_package_\myIndex} 
\fi

\def\myTitle{高等数学笔记}
\def\myAuthor{M.K.}
\def\myDateCover{\today}
\def\myDateForeword{\today}
\def\myForeword{前言}
\def\myForewordText{
    
    这是一个基于\LaTeX{}模板的高数笔记．
    内容参考了包括但不限于以下资料：
    \begin{itemize}
        \item 《高等数学》（同济大学出版社）第七版
        \item 《吉米多维奇高等数学习题精选精讲》（山东科学技术出版社）第二版
        \item 《高等数学作业集》（哈尔滨工业大学出版社）第一版
        \item \href{https://space.bilibili.com/1035929235}{B站UP主：一高数}的高数课程
        \item \href{https://space.bilibili.com/42428180}{B站UP主：考研竞赛数学}的试题讲解与高数课程
        \item \href{https://space.bilibili.com/391930545}{B站UP主：Maki的完美算数教室}的高数内容

    \end{itemize}

    封面来源于\href{https://www.pixiv.net/users/3952}{おにねこ}的作品\href{https://www.pixiv.net/artworks/136551599}{雨雫とプレアデス}
    如有侵权请联系我删除．
}
\def\mySubheading{副标题}


\begin{document}
% \input{../config/cover}
\else
\fi
\chapter{多元函数微分学}
\begin{exercise}
    设可微函数$z=z(x,y)$满足$x^2\dfrac{\partial z}{\partial x}+y^2\dfrac{\partial z}{\partial y}=2z^2$，又设$u=x,v=\dfrac{1}{y}-\dfrac{1}{x},w=\dfrac{1}{z}-\dfrac{1}{x}$，则对于函数$w=w(u,v)$，求$\dfrac{\partial w}{\partial u}\Big|_{\substack{u=2\\v=1}}$。
\end{exercise}
\begin{solution}
    已知 $u=x, v = \dfrac{1}{y} - \dfrac{1}{x}, w = \dfrac{1}{z} - \dfrac{1}{x}$．
    
    故$$x=u,\quad\dfrac{1}{y} = v + \dfrac{1}{u},\quad\dfrac{1}{z} = w + \dfrac{1}{u}$$
    
    对 $w=w(u,v)$ 关于 $u$ 求偏导：
    \[ \frac{\partial w}{\partial u} = \frac{\partial}{\partial u} \left( \frac{1}{z} - \frac{1}{u} \right) = -\frac{1}{z^2} \frac{\partial z}{\partial u} + \frac{1}{u^2} \]
    其中 $z$ 是 $x, y$ 的函数，而 $x, y$ 又是 $u, v$ 的函数．由复合函数求导法则：
    \[ \frac{\partial z}{\partial u} = \frac{\partial z}{\partial x} \frac{\partial x}{\partial u} + \frac{\partial z}{\partial y} \frac{\partial y}{\partial u} \]
    由 $x=u$ 知 $\dfrac{\partial x}{\partial u} = 1$．
    由 $\dfrac{1}{y} = v + \dfrac{1}{u}$ 知 $-\dfrac{1}{y^2} \dfrac{\partial y}{\partial u} = -\dfrac{1}{u^2}$，即 $\dfrac{\partial y}{\partial u} = \dfrac{y^2}{u^2} = \dfrac{y^2}{x^2}$．
    
    代入得：
    \[ \frac{\partial z}{\partial u} = \frac{\partial z}{\partial x} \cdot 1 + \frac{\partial z}{\partial y} \cdot \frac{y^2}{x^2} = \frac{1}{x^2} \left( x^2 \frac{\partial z}{\partial x} + y^2 \frac{\partial z}{\partial y} \right) \]
    已知原方程 $x^2 \dfrac{\partial z}{\partial x} + y^2 \dfrac{\partial z}{\partial y} = 2z^2$，则：
    \[ \frac{\partial z}{\partial u} = \frac{2z^2}{x^2} \]
    代入 $\dfrac{\partial w}{\partial u}$ 的表达式：
    \[ \frac{\partial w}{\partial u} = -\frac{1}{z^2} \cdot \frac{2z^2}{x^2} + \frac{1}{x^2} = -\frac{2}{x^2} + \frac{1}{x^2} = -\frac{1}{x^2} \]
    当 $u=2, v=1$ 时，$x=u=2$，故：
    \[ \frac{\partial w}{\partial u}\Big|_{\substack{u=2\\v=1}} = -\frac{1}{2^2} = -\frac{1}{4} \]
\end{solution}
\begin{exercise}
    设$z=f(x,y)$是区域$D=\{(x,y)|0\le x\le 1,0\le y\le 1\}$上连续可微的函数，且满足$f(0,0)=0$，且$\d z\Big|_{(0,0)}=3\d x+2\d y$，求极限$\displaystyle\lim_{x \to 0} \dfrac{\displaystyle \int_{0}^{x^2}\d t\displaystyle\int_{x}^{\sqrt{t}}f(t,u)\d u}{1-\sqrt[4]{1-x^4}}$
\end{exercise}
\begin{solution}
    设极限为 $L$．首先处理分母：当 $x \to 0$ 时，使用等价无穷小 $(1+x)^\alpha - 1 \sim \alpha x$，有：
    \[ 1 - \sqrt[4]{1-x^4} = 1 - (1 - x^4)^{1/4} \sim -\left[ \frac{1}{4} (-x^4) \right] = \frac{1}{4}x^4 \]
    于是：
    \[ L = \lim_{x \to 0} \frac{\displaystyle\int_{0}^{x^2} \d t \int_{x}^{\sqrt{t}} f(t,u) \d u}{\frac{1}{4}x^4} \]
    接下来处理分子，对分子先进行交换积分次序：
    \[ \int_{0}^{x^2} \d t \int_{x}^{\sqrt{t}} f(t,u) \d u = -\int_{0}^{x}\d u\int_{0}^{u^2} f(t,u) \d t \]
    容易发现整个分式是$\frac{0}{0}$型的，因此可以使用洛必达法则．对分子分母求导：
    \[L = \lim_{x \to 0^+}\frac{-\displaystyle\int_{0}^{x}\d u\int_{0}^{u^2} f(t,u) \d t}{\frac{1}{4}x^4}
        = -\lim_{x \to 0^+}\frac{\displaystyle\int_{0}^{x^2}f(t,x)\d t}{x^3} \]
    容易发现分子分母又是$\frac{0}{0}$型的，但如果继续使用洛必达法则，分子是变限积分，求导后积分仍然存在，不如使用定积分中值定理．\\
    由定积分中值定理，存在 $\xi \in (0, x^2)$ 使得：
    \[ \int_{0}^{x^2} f(t,x) \d t = f(\xi, x) \cdot x^2 \]
    因此：
    \[ L = -\lim_{x \to 0^+} \frac{f(\xi, x) \cdot x^2}{x^3} = -\lim_{x \to 0^+} \frac{f(\xi, x)}{x} \]
    当 $x \to 0^+$ 时，$\xi \to 0^+$，因此：
    \[ f(\xi, x) = f(0,0) + \frac{\partial f}{\partial x}(0,0) \cdot \xi + \frac{\partial f}{\partial y}(0,0) \cdot x + o(\sqrt{\xi^2 + x^2}) \]
    代入得：
    \[ L = -\lim_{x \to 0^+} \frac{f(0,0) + 3\xi + 2x + o(\sqrt{\xi^2 + x^2})}{x}
        = -2 + 3\lim_{x \to 0^+} \frac{\xi}{x}+\lim_{x\to 0^+}\frac{o(\sqrt{\xi^2 + x^2})}{x} \]
    由于 $0 < \xi < x^2$，所以$0 < \dfrac{\xi}{x} < x$，因此 $\lim\limits_{x \to 0^+} \dfrac{\xi}{x} = 0$．
    而 $$\lim_{x \to 0^+} \frac{\sqrt{\xi^2 + x^2}}{x} = \lim_{x \to 0^+} \frac{o(\sqrt{\xi^2 + x^2})}{\sqrt{\xi^2 + x^2}} \cdot \frac{\sqrt{\xi^2 + x^2}}{x} = 0$$
    因此：
    \[ L = -2 \]
    
\end{solution}
\begin{myRemark}
    关于多重变限积分的求导，事实上是莱布尼茨积分法则的推广，具体来说，如果 $F(x) = \displaystyle\int_{a(x)}^{b(x)} f(t,x) \d t$，则：
\[ F'(x) = f(b(x), x) \cdot b'(x) - f(a(x), x) \cdot a'(x) + \int_{a(x)}^{b(x)} \frac{\partial f(t,x)}{\partial x} \d t \]
    在本题中，分子中的积分是一个二重变限积分，因此需要对内外两层积分分别使用莱布尼茨法则进行求导．
\end{myRemark}
\ifx\allfiles\undefined
\end{document}
\fi